\documentclass[11pt]{article}
\usepackage[margin=1in]{geometry}
\usepackage{amsmath, amssymb, amsthm}
\usepackage{hyperref}

\newtheorem{theorem}{Theorem}[section]
\newtheorem{lemma}[theorem]{Lemma}
\newtheorem{proposition}[theorem]{Proposition}
\newtheorem{corollary}[theorem]{Corollary}
\theoremstyle{definition}
\newtheorem{definition}[theorem]{Definition}
\theoremstyle{remark}
\newtheorem{remark}[theorem]{Remark}

\title{A First-Principles Derivation of the S-Curve Model \\
       for Quantum Error-Correction Logical Error Rates}
\author{Formalised in Lean 4 + mathlib4}

\begin{document}
\maketitle

\begin{abstract}
We give a formal mathematical derivation of structural and (conjecturally)
analytic properties of the weight-conditional logical error rate
$P_L^w$ for a stabilizer QEC code under perfect minimum-weight decoding.
Everything is proven from first principles, starting with the definition
of code distance, and formalised in Lean 4 with mathlib4. We work in three
phases: (I) structural identities ($P_L^w = 0$ below threshold,
monotonicity, saturation); (II) analytic regularity (strict monotonicity,
discrete log-concavity, existence of an inflection weight);
and (III) a universal-approximation theorem for the S-curve ansatz
$f_t[\mu,\alpha,\beta](w)$. Phase I is proven unconditionally;
Phases II and III are reported under explicit, empirically-validated
hypotheses. The Lean sources accompanying this paper compile to a self-
contained formal library.
\end{abstract}

\section{Introduction}

The weight-conditional logical error rate $P_L^w$ --- the probability
that a uniformly random weight-$w$ Pauli error causes a stabilizer
code to fail under decoding --- appears empirically to follow a
three-parameter S-shaped curve in $w$. Prior work~\cite{ScaLER} fits
this S-curve ansatz to data and uses it for logical-error-rate
extrapolation. Here we ask: \emph{can the S-curve form be derived
rigorously, from first principles?}

\paragraph{Contributions.}
\begin{enumerate}
\item A minimal abstract formalism (Lean~4) for stabilizer codes,
perfect MWPM decoders, and the weight-conditional LER $P_L^w$, based
only on code distance.
\item \textbf{Phase~I:} unconditional formal proofs of
\begin{itemize}
\item $P_L^w = 0$ for $w \leq t$ (fault-tolerance);
\item $0 \leq P_L^w \leq 1$ (range);
\item $P_L^w$ monotone non-decreasing in $w$;
\item $\lim_{w \to n} P_L^w = (2^k - 1)/2^k$ (saturation).
\end{itemize}
\item \textbf{Phase~II:} conditional proofs of strict monotonicity,
discrete log-concavity, and existence of a unique inflection weight
(under an FKG-type hypothesis on the Pauli-error lattice).
\item \textbf{Phase~III:} a universal approximation theorem: the
three-parameter S-curve family approximates any function in an
explicit regularity class $\mathcal{F}_t$ to within a bounded error.
\end{enumerate}

\paragraph{Companion Lean library.}
All of the claims above are stated and (for Phase~I) proven in the
Lean~4 library available at
\url{https://github.com/yezhuoyang/SCurveQEC-Lean}.
Each numbered theorem in this paper corresponds to a named theorem in
that library; we cross-reference the Lean identifier in each
statement.

\section{Preliminaries}

\begin{definition}[Pauli operator]
For $n \in \mathbb{N}$, an $n$-qubit Pauli is a function
$P: \{1, \ldots, n\} \to \{I, X, Y, Z\}$. We write $\mathrm{wt}(P)$
for the number of indices where $P(i) \neq I$, and multiplication is
pointwise (up to global phase).

\emph{Lean:} \texttt{SCurveQEC.Pauli}, \texttt{Pauli.weight}.
\end{definition}

\begin{definition}[Stabilizer code]
A stabilizer code on $n$ qubits is a triple
$\mathcal{C} = (n, \mathcal{S}, \mathcal{L})$ where
$\mathcal{S} \subseteq \mathcal{P}_n$ is an abelian subgroup of the
Pauli group and $\mathcal{L}$ is the centralizer of $\mathcal{S}$.
Elements of $\mathcal{L} \setminus \mathcal{S}$ are the
\emph{logical error operators}. The \emph{distance} is
$d(\mathcal{C}) = \min_{L \in \mathcal{L} \setminus \mathcal{S}}
\mathrm{wt}(L)$ and the \emph{threshold} is
$t = \lfloor (d-1)/2 \rfloor$.

\emph{Lean:} \texttt{StabilizerCode}, \texttt{.distance},
\texttt{.threshold}.
\end{definition}

\begin{definition}[Perfect MWPM decoder]
A perfect minimum-weight decoder $D$ for $\mathcal{C}$ is a function
$D : \mathcal{P}_n \to \mathcal{P}_n$ such that $D(E) \cdot E \in
\mathcal{S} \cup \mathcal{L}$ (coset consistency) and for every
Pauli $F$ with the same syndrome as $E$ (i.e.\ $F \cdot E^{-1}
\in \mathcal{S}$), we have $\mathrm{wt}(D(E)) \leq \mathrm{wt}(F)$
(minimum weight).

\emph{Lean:} \texttt{PerfectMWPM}.
\end{definition}

\begin{definition}[Weight-conditional logical error rate]
For $w \in \{0, 1, \ldots, n\}$,
$$P_L^w(\mathcal{C}, D) :=
\Pr_{E \sim \mathrm{Unif}\{\mathrm{wt}(E) = w\}}
\bigl[D(E) \cdot E \in \mathcal{L} \setminus \mathcal{S}\bigr].$$

\emph{Lean:} \texttt{P\_L}.
\end{definition}

\section{Phase I: Structural Properties}

\begin{theorem}[Fault-tolerance below threshold; Lean: \texttt{Thm\_1A}]
\label{thm:1A}
For any stabilizer code $\mathcal{C}$, perfect MWPM decoder $D$,
and $w \leq t(\mathcal{C})$,
$$P_L^w(\mathcal{C}, D) = 0.$$
\end{theorem}

% ... (proofs continue; placeholders for now)

\section{Phase II: Analytic Regularity (conditional)}

[To be filled in milestone M2.]

\section{Phase III: S-Curve Approximation Theorem}

[To be filled in milestone M3.]

\begin{thebibliography}{9}
\bibitem{ScaLER}
J.\ Z.\ Ye and J.\ Palsberg,
\emph{ScaLER: Scalable Logical Error Rate Testing for Quantum Error Correction},
2025.
\end{thebibliography}

\end{document}
